% =======================================================
% 8. REQUISITI DI SISTEMA
% =======================================================
\section{Requisiti di Sistema}
\label{sec:requisiti-sistema}
Questa sezione esplicita i requisiti hardware, software e di rete assunti dall'architettura descritta nel documento, a cui i requisiti prestazionali del capitolo 5 -- Requisiti dell'\link{https://synergo-unit-unipd.github.io/Second-Brain/docs/PB/doc_tecnici/doc_esterni/analisi_dei_requisiti/analisi_dei_requisiti.pdf}{Analisi dei Requisiti} (in particolare R1-P-O, R3-P-O, R4-P-O) fanno riferimento senza definirli esplicitamente.

\subsection{Hardware}
Trattandosi di un'applicazione web in cui l'elaborazione AI è interamente delegata al servizio LLM esterno (Sezione~\ref{sec:arch-backend}), il carico computazionale lato client si limita al rendering dell'editor e dell'anteprima Markdown; non sono pertanto richieste risorse hardware dedicate oltre a quelle di un dispositivo desktop o laptop di fascia standard in commercio negli ultimi 3--4 anni.

\begin{itemize}
	\item \textbf{CPU}: processore dual-core con clock $\geq$ 1.6 GHz o equivalente, sufficiente a garantire i tempi di avvio (R3-P-O) e di aggiornamento dell'anteprima (R1-P-O, R4-P-O) dichiarati.
	\item \textbf{RAM}: 4 GB minimi consigliati per l'esecuzione del browser con una singola scheda dedicata all'applicazione, tenendo conto dell'overhead di CodeMirror\textsubscript{G} e del parsing Markdown lato client.
	\item \textbf{Storage}: nessun requisito significativo; il sistema non mantiene alcuna persistenza locale automatica oltre ai singoli file \texttt{.md} esplicitamente salvati dall'utente (R75-F-O), la cui dimensione è trascurabile rispetto alla capacità di un qualsiasi dispositivo moderno.
	\item \textbf{Lato server}: nessun requisito hardware specifico è imposto dal gruppo, poiché l'infrastruttura di calcolo per l'inferenza LLM è gestita e messa a disposizione dall'azienda proponente (R3-V-O); il Backend Second Brain richiede solo le risorse minime per l'esecuzione di un servizio FastAPI containerizzato (Sezione~\ref{sec:arch-deployment}).
\end{itemize}

Il degrado prestazionale sotto queste soglie non è stato misurato sistematicamente dal gruppo: i valori dichiarati in R1-P-O, R3-P-O e R4-P-O sono da intendersi validi per dispositivi conformi ai requisiti sopra indicati, mentre su hardware inferiore l'unica garanzia esplicita rimane l'assenza di blocchi o crash dell'editor (R10-Q-O).

\subsection{Software}
\begin{itemize}
	\item \textbf{Browser}: come vincolato da R5-V-O, il sistema deve essere eseguito su Chrome 151+, Edge 150+, Firefox 153+ o Safari 26+.
	\item \textbf{Sistema operativo}: nessun vincolo diretto oltre alla disponibilità di uno dei browser sopra indicati; l'applicazione è stata verificata su Windows, macOS e Linux.
	\item \textbf{Dipendenza da API browser-specifiche}: le funzionalità di salvataggio e apertura locale della nota (R75-F-O, R76-F-O) si appoggiano preferibilmente alla \textit{File System Access API} del browser (\texttt{showSaveFilePicker}, \texttt{showOpenFilePicker}, Sezione~\ref{sec:arch-frontend}), nativa su Chrome ed Edge. Su Firefox e Safari, dove questa API non è disponibile, \texttt{NoteServiceProxy} rileva l'assenza a runtime e ricade su un fallback equivalente basato su API standard (download tramite \texttt{Blob}/\texttt{<a download>} per il salvataggio, \texttt{<input type="file">} per l'apertura), garantendo così la compatibilità con l'intero parco browser richiesto da R5-V-O senza violare R13-Q-D (nessuna dipendenza da funzionalità non standard esposta al resto del sistema, dato che il contratto \texttt{NoteService} resta invariato in entrambi i casi).
\end{itemize}

\subsection{Network}
\begin{itemize}
	\item \textbf{Editing locale}: le funzionalità di scrittura, formattazione e cronologia delle modifiche (R1--R47) non richiedono connettività di rete, essendo interamente gestite lato client (R10-Q-O).
	\item \textbf{Operazioni AI}: richiedono una connessione di rete attiva tra Frontend e Backend e tra Backend e servizio LLM esterno; la comunicazione avviene su protocollo HTTPS (Sezione~\ref{sec:arch-backend}). Non è imposto un requisito esplicito di banda minima, ma il sistema applica un timeout massimo indicativo di 5 minuti per richiesta (R5-P-O), oltre il quale l'elaborazione viene interrotta con notifica di errore.
	\item \textbf{Persistenza remota (opzionale, non implementata)}: qualora in futuro venisse realizzato il modulo opzionale di persistenza remota (R4-V-Op), andrebbe definito un requisito di rete dedicato per il caricamento della lista delle note (R6-P-Op); allo stato attuale non è applicabile.
\end{itemize}

\FloatBarrier
